\documentclass[11pt]{article}
\usepackage{fontspec,amsmath,amssymb,amsthm,hyperref}
\usepackage[margin=1in]{geometry}
\setlength{\emergencystretch}{3em}
\title{Physical-time generator and Gibbs integration identity}
\author{LangevinLean mathematical exposition}
\date{}
\begin{document}\maketitle
\noindent Source module: \texttt{IPLA/Generator.lean}. Compile this proof document with XeLaTeX or LuaLaTeX. The accompanying HTML places the complete Lean declarations beside every proof.


For scaled state \(z\), physical IPLA time is \(N\) times the ordinary lifted Langevin time. The generator is consequently \(N^{-1}(\Delta-\nabla V_N\cdot\nabla)\).


\section{Identify the generator on a twice differentiable test function}


Define \(\mathcal A_N\varphi=N^{-1}\mathcal L_{V_N}\varphi\), where \(\mathcal L_V\) is the shared Langevin generator. For \(C^2\) test functions the shared coordinate definition equals the Laplacian minus the gradient inner product, hence

\[\mathcal A_N\varphi(z)=\frac1N\left(\Delta\varphi(z)-\langle\nabla V_N(z),\nabla\varphi(z)\rangle\right).\]

This is a direct identity of operators on the scaled state; no invariant-measure conclusion is built into its definition.



\noindent Lean declarations: \texttt{diffusionGenerator}, \texttt{diffusionGenerator\_eq}.

\section{Derive integrability and infinitesimal invariance}


For compactly supported \(C^2\) \(\varphi\), the shared Gibbs integration-by-parts theorem proves that \(\mathcal L_{V_N}\varphi\) is integrable and has Gibbs integral zero. Its potential hypotheses follow from differentiability and Lipschitzness of the joint gradient via the geometry module. Multiplication by the constant \(N^{-1}\) preserves integrability and gives

\[\int\mathcal A_N\varphi\,d\Pi_N=\frac1N\int\mathcal L_{V_N}\varphi\,d\Pi_N=0.\]

The theorem also returns the probability property, derived from positive strong convexity. Explicit integrability rules out a zero caused by the default integral value of a nonintegrable function. Full time-evolution stationarity is supplied separately by the shared Langevin construction.



\noindent Lean declarations: \texttt{gibbsTarget\_infinitesimalInvariant}.

\end{document}
